% manuscript_clean_2_background_excerpt.tex
% Self-contained excerpt of manuscript_clean_2.tex: front matter through Background.
% Bibliography entries used up to this point are embedded via filecontents (no external .bib).
% Compile: pdflatex -> bibtex -> pdflatex -> pdflatex

\begin{filecontents*}[overwrite]{\jobname.bib}
@Article{Abdulgader2018,
  author    = {Abdulgader, Fathia Sghayer and Eid, Rajeh and Daneshvar Rouyendegh, Babak},
  journal   = {Applied Computational Intelligence and Soft Computing},
  title     = {Development of Decision Support Model for Selecting a Maintenance Plan Using a Fuzzy {MCDM} Approach: A Theoretical Framework},
  year      = {2018},
  issn      = {1687-9732},
  month     = nov,
  pages     = {1--14},
  volume    = {2018},
  doi       = {10.1155/2018/9346945},
  publisher = {Wiley},
}

@Article{Arrieta2020,
  author  = {Barredo Arrieta, Alejandro and D{\'i}az-Rodr{\'i}guez, Natalia and Del Ser, Javier and Bennetot, Adrien and Tabik, Siham and Barbado, Alberto and Garcia, Salvador and Gil-Lopez, Sergio and Molina, Daniel and Benjamins, Richard and Chatila, Raja and Herrera, Francisco},
  journal = {Information Fusion},
  title   = {Explainable {Artificial Intelligence} ({XAI}): Concepts, Taxonomies, Opportunities and Challenges Toward Responsible {AI}},
  year    = {2020},
  issn    = {1566-2535},
  month   = jun,
  pages   = {82--115},
  volume  = {58},
  doi     = {10.1016/j.inffus.2019.12.012},
}

@Article{Borkotokey2019,
  author  = {Borkotokey, Surajit and Hazarika, Pankaj and Mesiar, Radko},
  journal = {International Journal of General Systems},
  title   = {Cooperative Games with Multiple Attributes},
  year    = {2019},
  issn    = {0308-1079},
  month   = sep,
  number  = {8},
  pages   = {825--842},
  volume  = {48},
  doi     = {10.1080/03081079.2019.1659792},
}

@Article{Chekry2024,
  author  = {Chekry, Abderrahman and Bakkas, Jamal and Hanine, Mohamed and Montero, Elizabeth Caro and {Garat de Marin}, Mirtha Silvana and Ashraf, Imran},
  journal = {SoftwareX},
  title   = {{PyDEMATEL}: A {Python}-based tool implementing {DEMATEL} and fuzzy {DEMATEL} methods for improved decision making},
  year    = {2024},
  doi     = {10.1016/j.softx.2024.101889},
}

@Article{Deliu2026,
  author  = {Deliu, Delia and Moldovan, Andreea-Loredana},
  journal = {International Journal of Managing Projects in Business},
  title   = {Human--{AI} Collaboration and Cognitive Flexibility in Industry 6.0: A Framework for Purposeful Decision-Making in Innovation and Project Management},
  year    = {2026},
  doi     = {10.1108/ijmpb-10-2025-0459},
}

@InBook{Godz2026,
  author    = {Godz, Mariia and Novoa-Hernández, Pavel and Pelta, David A.},
  pages     = {203--217},
  publisher = {Springer Nature Switzerland},
  title     = {A Cooperative and Context-Aware Approach for Employee Attrition Prevention},
  year      = {2026},
  isbn      = {9783032290007},
  booktitle = {Information Processing and Management of Uncertainty in Knowledge-Based Systems},
  doi       = {10.1007/978-3-032-29000-7_15},
  issn      = {1865-0937},
}

@Article{Grazioso2023,
  author  = {Grazioso, Stanislao and Caporaso, Teodorico and {di Gironimo}, Giuseppe},
  journal = {International Journal on Interactive Design and Manufacturing ({IJIDeM})},
  title   = {Decision support in engineering design: the {ELIGERE} open source software platform},
  year    = {2023},
  doi     = {10.1007/s12008-023-01568-2},
}

@Article{HerreraViedma2002,
  author  = {Herrera-Viedma, E. and Herrera, F. and Chiclana, F.},
  journal = {{IEEE} Transactions on Systems, Man, and Cybernetics---Part A: Systems and Humans},
  title   = {A Consensus Model for Multiperson Decision Making with Different Preference Structures},
  year    = {2002},
  number  = {3},
  pages   = {394--402},
  volume  = {32},
  doi     = {10.1109/TSMCA.2002.802821},
}

@Article{Kacprzyk2022,
  author    = {Kacprzyk, Janusz and Sirbiladze, Gia and Tsulaia, Giorgi},
  journal   = {International Journal of Information Technology \& Decision Making},
  title     = {Associated Fuzzy Probabilities in {MADM} with Interacting Attributes: Application in Multi-objective Facility Location Selection Problem},
  year      = {2022},
  number    = {4},
  pages     = {1155--1188},
  volume    = {21},
  doi       = {10.1142/s0219622022500146},
  publisher = {World Scientific},
}

@InCollection{Lamata2020,
  author    = {Maria T. Lamata and David A. Pelta and Jos{\'{e}} Luis Verdegay},
  booktitle = {Fuzzy Approaches for Soft Computing and Approximate Reasoning: Theories and Applications},
  publisher = {Springer International Publishing},
  title     = {The Role of the Context in Decision and Optimization Problems},
  year      = {2020},
  month     = {oct},
  pages     = {75--84},
  doi       = {10.1007/978-3-030-54341-9_7},
}

@Article{Li2024,
  author    = {Li, Mingwei and Sun, Hao and Huang, Yingying and Chen, Hailin},
  journal   = {Autonomous Intelligent Systems},
  title     = {Shapley Value: From Cooperative Game to Explainable Artificial Intelligence},
  year      = {2024},
  number    = {1},
  pages     = {2},
  volume    = {4},
  doi       = {10.1007/s43684-023-00060-8},
  publisher = {Springer},
}

@InProceedings{Lundberg2017,
  author    = {Lundberg, Scott M. and Lee, Su-In},
  booktitle = {Advances in Neural Information Processing Systems},
  title     = {A Unified Approach to Interpreting Model Predictions},
  year      = {2017},
  pages     = {4765--4774},
  publisher = {Curran Associates},
  volume    = {30},
  url       = {https://proceedings.neurips.cc/paper/2017/hash/8a20a8621978632d76c43dfd28b67767-Abstract.html},
}

@Article{Lundberg2020,
  author  = {Lundberg, Scott M. and Erion, Gabriel and Chen, Hugh and DeGrave, Alex and Prutkin, Jordan M. and Nair, Bala and Katz, Ronit and Himmelfarb, Jonathan and Bansal, Nisha and Lee, Su-In},
  journal = {Nature Machine Intelligence},
  title   = {From Local Explanations to Global Understanding with Explainable {AI} for Trees},
  year    = {2020},
  number  = {1},
  pages   = {56--67},
  volume  = {2},
  doi     = {10.1038/s42256-019-0138-9},
}

@Article{NavarroGalera2025,
  author    = {Navarro-Galera, Andrés and Lara-Rubio, Juan and Novoa-Hernández, Pavel and Cruz Corona, Carlos A.},
  journal   = {Computational Economics},
  title     = {Using Decision Trees to Predict Insolvency in Spanish SMEs: Is Early Warning Possible?},
  year      = {2025},
  issn      = {1572-9974},
  month     = mar,
  number    = {1},
  pages     = {91–116},
  volume    = {65},
  doi       = {10.1007/s10614-024-10586-5},
  publisher = {Springer Science and Business Media LLC},
  url       = {http://dx.doi.org/10.1007/s10614-024-10586-5},
}

@InProceedings{NovoaHernandez2026,
  author    = {Novoa-Hernández, Pavel and Pelta, David A. and Godz, Mariia and Verdegay, José L. and Buendia-Carrillo, Dionisio},
  booktitle = {2026 IEEE Conference on Artificial Intelligence (CAI)},
  title     = {CADEMAS – A Framework for Cooperative Automated Decision-Making Systems},
  year      = {2026},
  pages     = {756-761},
  doi       = {10.1109/CAI68641.2026.11536392},
}

@Article{Ozkurt2025,
  author  = {\"Ozkurt, Cem and Canay, \"Ozkan and Tun{\c{c}}, Ey{\"u}p Altu{\u{g}} and Aydin, Elif and Velio{\u{g}}lu, Beyza S{\i}la},
  journal = {Baltic Journal of Modern Computing},
  title   = {Renewable Energy Source Ranking and Analysis Using Fuzzy {MCDM}, {ML}, and {XAI} Techniques},
  year    = {2025},
  number  = {3},
  pages   = {656--679},
  volume  = {13},
  doi     = {10.22364/bjmc.2025.13.3.06},
}

@Article{PerezCanedo2023,
  author    = {Boris P{\'{e}}rez-Ca{\~{n}}edo and Cynthia Porras and David A. Pelta and Jos{\'{e}} Luis Verdegay},
  journal   = {{IEEE} Transactions on Fuzzy Systems},
  title     = {Modeling Contexts as Fuzzy Propositions in Optimization Problems},
  year      = {2023},
  month     = {may},
  number    = {5},
  pages     = {1474--1483},
  volume    = {31},
  doi       = {10.1109/tfuzz.2022.3203786},
  publisher = {Institute of Electrical and Electronics Engineers ({IEEE})},
}

@Article{PerezDominguez2024,
  author  = {P{\'e}rez-Dom{\'\i}nguez, Luis and Callejas-Cuervo, M. and Garc{\'\i}a-Luna, F.},
  journal = {IEEE Access},
  title   = {Intuitionistic Fuzzy Dimensional Analysis in Multi-Criteria Decision-Making: A Computational Approach},
  year    = {2024},
  pages   = {1--15},
  volume  = {12},
  doi     = {10.1109/ACCESS.2024.3422616},
}

@Article{Roszkowska2024,
  author  = {Roszkowska, Ewa and Jefma{\'n}ski, Bart{\l}omiej and Dudek, Andrzej and Kusterka-Jefma{\'n}ska, Marta},
  journal = {SoftwareX},
  title   = {{IFMCDM}: An {R} package for intuitionistic fuzzy multi-criteria decision making methods},
  year    = {2024},
  pages   = {101676},
  volume  = {26},
  doi     = {10.1016/j.softx.2024.101721},
}

@Article{Sahraei2022,
  author  = {Sahraei, R. and Kanani-Sadat, Y. and Homayouni, Saeid and Safari, A. and Oubennaceur, Khalid and Chokmani, K.},
  journal = {Journal of Flood Risk Management},
  title   = {A Novel Hybrid {GIS}-Based Multi-Criteria Decision-Making Approach for Flood Susceptibility Analysis in Large Ungauged Watersheds},
  year    = {2022},
  number  = {2},
  pages   = {e12871},
  volume  = {16},
  doi     = {10.1111/jfr3.12871},
}

@Article{Salo1995,
  author  = {Salo, Ahti A.},
  journal = {European Journal of Operational Research},
  title   = {Interactive Decision Aiding for Group Decision Support},
  year    = {1995},
  number  = {1},
  pages   = {134--149},
  volume  = {84},
  doi     = {10.1016/0377-2217(94)00322-1},
}

@Article{Sun2024,
  author  = {Sun, Chuanneng and Huang, Songjun and Pompili, Dario},
  journal = {AI},
  title   = {{LLM}-based Multi-Agent Reinforcement Learning: Current and Future Directions},
  year    = {2024},
  number  = {3},
  pages   = {1234--1252},
  volume  = {5},
  doi     = {10.3390/ai5030065},
}

@Article{Tariq2024a,
  author  = {Tariq, Shahroz and Chhetri, Mohan Baruwal and Nepal, Surya and Paris, C{\'e}cile},
  journal = {Expert Systems with Applications},
  title   = {{A2C}: A Modular Multi-stage Collaborative Decision Framework for Human-{AI} Teams},
  year    = {2024},
  pages   = {122334},
  volume  = {238},
  doi     = {10.1016/j.eswa.2023.122334},
}

@Article{Vahidov2004,
  author  = {Vahidov, Rustam and Fazlollahi, Bijan},
  journal = {Information \& Management},
  title   = {Pluralistic Multi-agent Decision Support System: A Framework and an Empirical Test},
  year    = {2004},
  issn    = {0378-7206},
  number  = {7},
  pages   = {883--898},
  volume  = {41},
  doi     = {10.1016/j.im.2003.08.017},
}

@Article{Vazquez2026,
  author        = {Vazquez-S{\'{a}}nchez, Angel Alberto and Buend{\'{i}}a-Carrillo, Dionisio and Corona, Carlos Cruz},
  journal       = {Mathematics},
  title         = {A Cooperative Automated Decision-Making Strategy with Fuzzy Contextual Modulation for {SME} Insolvency Prediction},
  year          = {2026},
  number        = {19},
  volume        = {14},
  article-number = {3516},
  doi           = {10.3390/math14193516},
  issn          = {2227-7390},
  url           = {https://www.mdpi.com/2227-7390/14/19/3516},
}

@Article{Wieckowski2022,
  author    = {Wi{\k{e}}ckowski, Jakub and Kizielewicz, Bart{\l}omiej and Sa{\l}abun, Wojciech},
  journal   = {SoftwareX},
  title     = {{pyFDM}: A {Python} library for uncertainty decision analysis methods},
  year      = {2022},
  issn      = {2352-7110},
  month     = dec,
  pages     = {101271},
  volume    = {20},
  doi       = {10.1016/j.softx.2022.101271},
  publisher = {Elsevier BV},
}

@Article{Wieckowski2023,
  author    = {Wi{\k{e}}ckowski, Jakub and Kizielewicz, Bart{\l}omiej and Sa{\l}abun, Wojciech},
  journal   = {SoftwareX},
  title     = {Handling decision-making in intuitionistic fuzzy environment: {PyIFDM} package},
  year      = {2023},
  issn      = {2352-7110},
  month     = may,
  pages     = {101344},
  volume    = {22},
  doi       = {10.1016/j.softx.2023.101344},
  publisher = {Elsevier BV},
}

@Article{Wieckowski2024,
  author  = {Wi{\k{e}}ckowski, Jakub and Sa{\l}abun, Wojciech},
  journal = {SoftwareX},
  title   = {{MakeDecision}: Online system for the graphical design of decision-making models in crisp and fuzzy environments},
  year    = {2024},
  doi     = {10.1016/j.softx.2024.101658},
}

@article{bui1999agent,
title = {An agent-based framework for building decision support systems},
journal = {Decision Support Systems},
volume = {25},
number = {3},
pages = {225-237},
year = {1999},
issn = {0167-9236},
doi = {10.1016/S0167-9236(99)00008-1},
author = {Tung Bui and Jintae Lee},
}

@Article{novoa2026aggregation,
  author        = {Novoa-Hern{\'{a}}ndez, Pavel and Godz, Mariia and Pelta, David A.},
  journal       = {Mathematics},
  title         = {Aggregation Semantics for Prioritization in Cooperative Automated Decision-Making},
  year          = {2026},
  number        = {18},
  volume        = {14},
  article-number = {3397},
  doi           = {10.3390/math14183397},
  issn          = {2227-7390},
  url           = {https://www.mdpi.com/2227-7390/14/18/3397},
}
\end{filecontents*}

\documentclass{interact}

\usepackage[T1]{fontenc}
\usepackage[utf8]{inputenc}
\usepackage{epstopdf} 
\usepackage{graphicx}
\usepackage{booktabs}
\usepackage{array}
\usepackage{tabularx}
\usepackage{amsmath}
\usepackage{amssymb}
\usepackage{url}
\usepackage{xcolor}
\usepackage{mdframed}
\usepackage{setspace}
\usepackage{tikz}
\usetikzlibrary{arrows.meta,positioning,fit,calc}

% Slightly more open line spacing for readability
\setstretch{1.25}

% Light gray callout for NLG quotes: box matches default quote width
\definecolor{quoteboxbg}{gray}{0.94}
\mdfdefinestyle{quotebox}{%
 backgroundcolor=quoteboxbg,
 hidealllines=true,
 leftmargin=\leftmargin,
 rightmargin=\leftmargin,
 innertopmargin=10pt,
 innerbottommargin=10pt,
 innerleftmargin=10pt,
 innerrightmargin=10pt,
 skipabove=\topsep,
 skipbelow=\topsep,
}
\renewenvironment{quote}{%
 \begin{mdframed}[style=quotebox]%
}{%
 \end{mdframed}%
}

% Float placement: prefer top of page, keep floats near their callouts,
% and flush deferred floats before major section breaks when needed.
\usepackage{flafter}
\usepackage{placeins}
\renewcommand{\topfraction}{0.95}
\renewcommand{\bottomfraction}{0.8}
\renewcommand{\textfraction}{0.05}
\renewcommand{\floatpagefraction}{0.8}
\setcounter{topnumber}{3}
\setcounter{bottomnumber}{2}
\setcounter{totalnumber}{4}
\setlength{\floatsep}{4pt plus 1pt minus 1pt}
\setlength{\textfloatsep}{8pt plus 2pt minus 2pt}
\setlength{\intextsep}{8pt plus 2pt minus 2pt}
\raggedbottom

% APA author-year (Taylor & Francis Interact + APA)
\usepackage[natbibapa,nodoi]{apacite}
\setlength\bibhang{12pt}
\renewcommand\bibliographytypesize{\fontsize{10}{12}\selectfont}

\usepackage{hyperref}

\begin{document}

\articletype{Software Tool Article} 

\title{CADEMAS-ML: A Software Tool for Context-Aware Cooperative and Explainable Automated Decision-Making}

\author{
\name{Pavel Novoa-Hern\'{a}ndez\textsuperscript{a}\thanks{CONTACT Pavel Novoa-Hern\'{a}ndez. Email: pnovoahe@ull.edu.es} and David A. Pelta\textsuperscript{b}}
\affil{\textsuperscript{a}Dept. Computer Engineering and Systems, Universidad de La Laguna, San Crist\'{o}bal de La Laguna, Spain; \textsuperscript{b}Dept. Computer Science and Artificial Intelligence, Universidad de Granada, Granada, Spain}
}

\maketitle

\begin{abstract}
Complex organizational decisions often require several actors to combine their assessments while accounting for the surrounding operational and strategic context. The Cooperative Automated Decision-Making Systems (CADEMAS) framework was recently proposed to represent this form of cooperative and context-aware decision-making. Previous CADEMAS instantiations, however, cover only specific configurations of a framework that admits substantially broader combinations of automated decision components, contextual models, and aggregation mechanisms. This flexibility is valuable in practice, but it also makes the framework difficult to operationalize without dedicated software support. This paper presents CADEMAS-ML, a software tool that implements a configurable class of CADEMAS systems in which automated decision-making components are machine-learning models and contextual knowledge is represented through fuzzy logic. The tool integrates and weights multiple predictive models, supports configurable fuzzy contexts and alternative aggregation operators, and makes the trade-off between predictive risk and contextual suitability explicit. It also incorporates explainability and robustness analysis to help users inspect and assess the resulting decisions. An employee-attrition case study illustrates how the same predictive evidence can lead to different prioritizations under alternative contextual scenarios.
\end{abstract}

\begin{keywords}
automated decision-making systems; fuzzy decision making; context-aware decision support; machine learning; prioritization; cooperative systems
\end{keywords}


\section{Introduction}\label{sec:intro}

Decision-making in real-world organizations is rarely a matter of applying a single analytical model to a well-defined set of inputs. Complex decisions often involve multiple stakeholders, organizational units, or analytical components that contribute different pieces of evidence to the same problem. These contributions may reflect different objectives, information sources, areas of expertise, or levels of reliability, and therefore need to be coordinated before they can support a coherent course of action.

This perspective becomes particularly important when automated decision-making is used to support prediction, classification, or risk assessment. An automated decision-making (ADM) component can produce a probability, rating, recommendation, or risk score, but that output alone does not determine how the result should be acted upon. The appropriate course of action may also depend on organizational priorities, available resources, departmental objectives, regulatory constraints, or the strategic conditions surrounding the decision. In short, decisions are taken in a context, not in a vacuum.

Different ADM components may also provide complementary views of the same cases, while the final decision may need to reflect contextual knowledge provided by domain experts or organizational policies.

These considerations motivated the Cooperative Automated Decision-Making Systems (CADEMAS) framework proposed by \citet{NovoaHernandez2026}. CADEMAS coordinates heterogeneous ADM components while incorporating the context in which their outputs are interpreted. Its central idea is to keep the generation and integration of decision evidence separate from the contextual evaluation of that evidence, so that a collective assessment can later be combined with organizational priorities, policies, or scenario-dependent considerations. Section~\ref{sec:cademas} formalizes this architecture.

The framework builds on cooperative and multi-agent decision-support systems, in which heterogeneous components contribute to a common decision process \citep{Vahidov2004,bui1999agent}, and on multiperson and group decision-support research, where heterogeneous information and preferences are traditionally combined through aggregation, negotiation, and consensus-seeking mechanisms \citep{Salo1995,HerreraViedma2002}. Similar principles extend naturally to automated settings, where the cooperating components are software-based decision units rather than human decision makers.

A particularly relevant scenario is prioritization under a limited intervention budget: given a set of entities, which ones should receive attention first? Examples include attrition-risk prevention \cite{Godz2026}, financial rescue for small and medium-sized companies \cite{NavarroGalera2025}, and student-dropout prevention (BUSCAR REF).

In such problems, CADEMAS combines the assessments of multiple ADMs with the suitability of each entity for a given decision context. For instance, under an economic-crisis policy one may prioritize companies with many employees, whereas under a growth policy one may favor those with high expected impact. This connection is developed formally in the full manuscript. Fuzzy rules offer a convenient way to represent such contexts \citep{Lamata2020,PerezCanedo2023}, combining quantitative information with expert knowledge and linguistic assessments.

Users may also need to understand why a particular entity receives a given prediction, especially when predictions contribute to decisions with organizational or human consequences \citep{Arrieta2020}. Local feature attributions can help inspect the factors associated with individual predictions \citep{Lundberg2017,Lundberg2020}.

In addition, the ordering of cases may change when the contextual definition, the aggregation mechanism, or the balance assigned to predictive and contextual information is revised, even if the underlying models remain unchanged. Robustness analysis reveals which priorities are stable and which cases are particularly sensitive to these policy choices \citep{Arrieta2020}.

To the best of our knowledge, no integrated tool currently supports this complete prioritization workflow.

This paper presents CADEMAS-ML to address that gap. The tool implements a configurable class of CADEMAS systems for prioritization problems, with machine-learning ADM units and fuzzy contextual knowledge. Beyond the core CADEMAS capabilities, it incorporates case-level explanation and robustness analysis as natural extensions of the workflow.


The remainder of this excerpt reviews the CADEMAS framework and related work on cooperative and fuzzy decision-support systems (Section~\ref{sec:background}).


\section{Background and related works}\label{sec:background}

This section sets out the conceptual foundations on which CADEMAS-ML builds. We first revisit the CADEMAS framework in its general form (Section~\ref{sec:cademas}). We then situate the tool among related cooperative-decision frameworks and application-oriented decision-support systems, clarifying the gap it is designed to fill (Section~\ref{sec:related}).

\subsection{CADEMAS framework}\label{sec:cademas}

CADEMAS \citep{NovoaHernandez2026} provides a general conceptual and architectural framework for structuring cooperation among heterogeneous Automated Decision-Making (ADM) units. Organizational decisions may draw on classifiers, optimization or simulation models, rule-based systems, software agents, or institutional procedures, whose assumptions, data requirements, output types, and reliability may differ substantially. Rather than prescribing a particular learning method or coordination protocol, CADEMAS provides a common structure through which these heterogeneous components can contribute to a shared decision process and their outputs can subsequently be interpreted in context.

A central premise of the framework is that an ADM unit need not produce a complete operational decision on its own. Instead, it may contribute a \emph{partial decision signal}, such as a score, prediction, rule activation, explanation, or suggested action. The semantics of these signals need not be identical across units, as cooperation does not require the underlying ADM mechanisms to be homogeneous. Their outputs are therefore brought together through an integration layer and subsequently considered in relation to the context in which the decision is made. Accordingly, a CADEMAS instance is represented by the tuple
\begin{equation*}
S=\langle \mathcal{IM},\,A,\,\mathcal{OM},\,\mathcal{C},\,\mathcal{T}\rangle,
\end{equation*}
where $\mathcal{IM}$ is the Input Manager, $A$ is the set of ADM units, $\mathcal{OM}$ is the Decision Integration Layer, $\mathcal{C}$ is the Contextual Modulator, and $\mathcal{T}$ is the universe of admissible data types and representations.

The Input Manager $\mathcal{IM}$ is responsible for routing, transforming, and validating the information entering the architecture. In particular, it distinguishes training or calibration data, which are used offline to construct or update the internal models of the ADM units, from decision-case data provided when a new case is evaluated. The ADM units are denoted by
\begin{equation*}
A=\{ADM_1,\ldots,ADM_n\}.
\end{equation*}
Each $ADM_k$ receives a data representation $r_k\in\mathcal{T}$ and produces an output $o_k\in\mathcal{T}$. Depending on the application, cooperation may be \emph{indirect}, with each unit contributing an independent signal, or \emph{direct}, with units exchanging intermediate representations. CADEMAS accommodates both forms while leaving the specific communication protocol to the particular application.

The outputs generated by the ADM units form the heterogeneous collection $O=\{o_1,\ldots,o_n\}$. The Decision Integration Layer $\mathcal{OM}$ brings these outputs together and synthesizes them into a unified decision object $D$ through the mapping
\begin{equation*}
\Phi:O\rightarrow D.
\end{equation*}
This integration establishes a common representation for the information contributed by the different ADM units and provides the basis for its subsequent consideration in context. The resulting $D$ is a \emph{raw} decision proposal, since the outputs have been harmonized but have not yet been adjusted to the conditions under which the decision is to be used.

The Contextual Modulator then incorporates the relevant contextual information into this integrated proposal. Formally,
\begin{equation*}
\Gamma:D\times\mathcal{C}\rightarrow D,\qquad
D^{\mathrm{final}}=\Gamma(D^{\mathrm{raw}},\mathcal{C}),
\end{equation*}
where $\mathcal{C}$ represents the contextual conditions relevant to the decision. These may include regulatory requirements, organizational policies, resource constraints, or other scenario-dependent considerations. Depending on the application, contextual information may be represented through rules, utility functions, constraints, fuzzy propositions, or other complementary mechanisms. The resulting $D^{\mathrm{final}}$ therefore represents the decision after the information provided by the ADM units has been integrated and considered in its relevant context.

\subsection{Related frameworks and decision-support applications}\label{sec:related}

Recent research has increasingly recognized that effective decision support requires more than accurate isolated models. In many application domains, decisions emerge from the interaction between automated units, expert knowledge, and contextual considerations that may change over time. For clarity, it is useful to distinguish two related but different bodies of work: general cooperative-decision frameworks, which address how multiple decision entities interact, and application-oriented decision-support systems, which implement concrete methods for a domain or family of decision problems. 

At the framework level, cooperative game theory offers one formal perspective on how decision-making entities with heterogeneous roles or attributes may combine their contributions, including settings in which each agent is characterized by multiple attributes rather than a single payoff value \citep{Borkotokey2019}. CADEMAS does not adopt this game-theoretic formulation, but shares its concern with making cooperation among heterogeneous entities explicit. At the operational level, \citet{Tariq2024a} propose an AI-to-collaborator (A2C) framework in which humans and AI systems share decision authority through explicit interaction mechanisms. Related work on hybrid intelligence emphasizes that authority distribution and human interpretive capacity should be modeled separately, because the level of AI integration and the cognitive flexibility of human decision-makers jointly shape decision quality and accountability \citep{Deliu2026}. A neighboring line of work examines cooperation among multiple AI agents, including LLM-based multi-agent reinforcement learning \citep{Sun2024}. These contributions provide relevant points of comparison because they address cooperation among decision entities. However, their main focus is not the reusable implementation of a CADEMAS-like workflow in which heterogeneous ADM outputs are explicitly separated from context modulation and final prioritization.

A related theoretical thread connects cooperative game theory to case-level explanation. The Shapley value distributes a joint outcome among players according to their averaged marginal contributions \citep{Borkotokey2019} and has been adopted in explainable artificial intelligence as a principled attribution mechanism. SHapley Additive exPlanations (SHAP) \citep{Lundberg2017} treat model features as players and decompose each prediction into additive feature contributions with axiomatic guarantees; recent surveys document the breadth of this cooperative-to-XAI bridge \citep{Li2024}. This literature motivates the inclusion of an explicit explanation module in CADEMAS-ML. The current implementation does not compute Shapley values over features: because the uploaded H2O MOJO models are treated as black boxes and may include stacked ensembles, the tool uses a model-agnostic perturbation procedure instead. The purpose is nonetheless the same as in Shapley-inspired XAI: to make the local drivers of a recommendation inspectable alongside the cooperative and contextual scores.

At the application and decision-support level, substantial progress has been made in fuzzy decision support, multi-criteria decision making, and hybrid fuzzy-ML workflows. Applications such as \textsc{MakeDecision} \citep{Wieckowski2024} and the fuzzy DEMATEL-based tool of \citet{Chekry2024} emphasize transparency and traceability by exposing intermediate calculations to users. A related stream has explored hybrid pipelines that combine expert judgment, causal weighting, and predictive components in domain-specific settings \citep{Abdulgader2018,Sahraei2022,Ozkurt2025}. Recent work has also extended the fuzzy toolbox through aggregation operators that account for interactions among decision attributes, capturing dependency structures that simple weighted averages cannot represent \citep{Kacprzyk2022}. The importance of explainability in AI-based decision support has been broadly recognized: any automated system whose outputs must be justified or audited requires transparent mechanisms to convey its reasoning to stakeholders \citep{Arrieta2020}, a requirement increasingly addressed through XAI techniques, including Shapley-value attribution \citep{Lundberg2017,Li2024}, in hybrid ranking and prediction pipelines \citep{Ozkurt2025}. These contributions demonstrate the value of combining data-driven evidence with contextual or expert-defined criteria, but most are tailored to particular methods or domains rather than organized around the CADEMAS decomposition of ADM units, integration, context, and final prioritization.

Open-source fuzzy and MCDM libraries, including \textsc{pyFDM} \citep{Wieckowski2022}, \textsc{PyIFDM} \citep{Wieckowski2023}, the \textsc{IFMCDM} package of \citet{Roszkowska2024}, the \textsc{Eligere} platform \citep{Grazioso2023}, and the intuitionistic fuzzy dimensional-analysis approach of \citet{PerezDominguez2024}, are also relevant. They improve reproducibility and provide reusable computational building blocks.

\begin{table}[!tb]
\tbl{Comparison of representative application-oriented decision-support works related to CADEMAS-ML.}
{\footnotesize
\begin{tabularx}{\textwidth}{>{\raggedright\arraybackslash}p{1.85cm} >{\raggedright\arraybackslash}X >{\centering\arraybackslash}p{1.35cm} >{\centering\arraybackslash}p{1.35cm} >{\centering\arraybackslash}p{1.35cm} >{\raggedright\arraybackslash}X}
 \toprule
 \textbf{Work} & \textbf{Decision-support focus} & \textbf{Multiple ADM units} & \textbf{Explicit context layer} & \textbf{Case-level XAI} & \textbf{Relation to CADEMAS-ML} \\
 \midrule
 \citet{Tariq2024a} & Human--AI collaborative decision framework with automated, augmented, and collaborative modes & Partly & Partly & No & Strong framework-level precedent, but not a configurable fuzzy/ML prioritization workflow \\
 \citet{Wieckowski2024} & Auditable fuzzy MCDM through a web-based tool & No & Partly & No & Strong transparency, but no cooperative ADM integration layer \\
 \citet{Chekry2024} & Open-source fuzzy DEMATEL decision-support tool & No & Partly & No & Traceable fuzzy causal analysis, but method-specific rather than CADEMAS-like \\
 \citet{Ozkurt2025} & Hybrid fuzzy MCDM, ML, and XAI for transparent ranking and prediction & Partly & Partly & Partly & Combines ranking, prediction, and XAI, but lacks a general cooperative ADM architecture \\
 \citet{Sahraei2022} & Fuzzy DEMATEL/BWM/ANP pipeline for spatial risk prioritization & No & Partly & No & Explicit causal weighting under uncertainty, but application-specific \\
 \midrule
 \citet{Godz2026} & Employee-attrition prevention using CADEMAS with departmental ML models & Yes & Yes & No & Direct CADEMAS application that motivates the implemented workflow \\
 \citet{novoa2026aggregation} & Aggregation semantics for CADEMAS prioritization under alternative operators & Yes & Yes & No & Analyzes how linear, geometric, and minimum rules preserve or compensate contextual requirements \\
 \citet{Vazquez2026} & Cooperative ADM with fuzzy contextual modulation for SME insolvency prediction & Yes & Yes & Yes & Extends CADEMAS-style cooperation to survival models with contextual modulation and XAI support \\
 \midrule
CADEMAS-ML (this work) & Configurable prioritization with ML-based ADM units, fuzzy context scenarios, and local case explanation & Yes & Yes & Yes & Operationalizes CADEMAS for reusable, explainable ML-driven decision support \\
 \bottomrule
	\end{tabularx}}
\label{tab:related}
\end{table}

Table~\ref{tab:related} summarizes representative application-oriented contributions in these areas. Three observations emerge from this comparison. First, published systems already provide strong ingredients for transparent decision support, including human--AI collaboration, auditable fuzzy MCDM, hybrid ranking pipelines, and reusable software libraries \citep{Tariq2024a,Wieckowski2024,Ozkurt2025,Wieckowski2022}. Second, these ingredients are usually developed in isolation. Third, recent CADEMAS applications already combine cooperative multi-ADM integration with an explicit contextual layer \citep{Godz2026,novoa2026aggregation,Vazquez2026}, yet they remain domain-specific configurations rather than a reusable software workflow with case-level explanation. What remains missing is an operational implementation that keeps these capabilities configurable and auditable together.

CADEMAS-ML is designed to address this gap for the predictive-model setting. Building on the framework and its attrition instantiation discussed above, it provides a configurable software implementation in which ML-based ADM outputs, contextual reasoning, and case-level explanation coexist as explicit and auditable components of the decision-making process. The source code is publicly available under an MIT license; the main contribution is the operationalization of the CADEMAS decomposition in a reusable decision-support tool.

\bibliographystyle{apacite}
\bibliography{\jobname}

\end{document}
